% Common-alphabet auxiliary route for SCP10 Theorem 5.5.
% Source: eq:2d:closure-intersection, eq:2d:closure-inv, eq:2d:close-in-in;
% coefficient duality: Lemma 4.6. The full theorem has a separate owner.
\subsection{Bond support and common-alphabet closure reconstruction}
\label{sec:peps_four_cut_common_reconstruction}

Throughout this subsection $G$ and $V$ are finite, and the torus periods
are positive. The site tensors $A^0_v$ are the matched averaging tensors
of Definition~\ref{def:peps_four_cut_common_matched_action}.
Let $\mathcal B=X\times\{\mathrm h,\mathrm v\}$ denote the actual bonds.

\begin{definition}[Bond labels, product vectors, and trace-dual coefficients]%
    \label{def:peps_four_cut_common_bond_coefficients}
    \lean{TNLean.PEPS.TorusBondLabels,
        TNLean.PEPS.torusRepresentationBondProduct,
        TNLean.PEPS.torusBondCoefficientExtraction,
        TNLean.PEPS.torusBondLabelConfigEquiv,
        TNLean.PEPS.card_torusBondLabels_two}
    \leanok
    \uses{def:peps_four_cut_common_matched_action,
        def:peps_torus_physical_bond_regrouping,
        thm:peps_torus_projector_extraction}
    A bond labelling is $p=(p^{\mathrm h},p^{\mathrm v})\in G^X\times G^X$.
    The bijection to $G^{\mathcal B}$ sends $p$ to
    $(v,d)\mapsto p^d_v$; its inverse restricts to the two directions.
    For any finite label alphabet $G$, there are $|G|^8$ such labellings
    on the two-by-two torus. Define the bare product vector
    \begin{align}
        B_p(\sigma) & =\prod_{v\in X}
        U^{\mathrm h}_v(p^{\mathrm h}_v)(l_{v+e},r_v)
        U^{\mathrm v}_v(p^{\mathrm v}_v)(t_v,b_{v+n}).
        \label{eq:peps_common_bond_product}
    \end{align}
    For each bond $f$, let $U_f$ denote its representation and put
    $\delta_{f,k}(M)=\tr(U_f(k^{-1})M\Delta_{U_f})$, using the trace-dual
    operator of Theorem~\ref{thm:peps_torus_projector_extraction}.
    The linear coefficient functional is
    \begin{align}
        E_p(\psi) & =\sum_{\beta:\mathcal B\to V^2}
        (\mathcal R\psi)(\beta)
        \prod_{f\in\mathcal B}\delta_{f,p_f}
        (E_{\beta_f^1,\beta_f^2}).
        \label{eq:peps_common_bond_extraction}
    \end{align}
    For unitary representations, $U_f(k^{-1})=U_f(k)^\dagger$,
    recovering the pairing in \cite[Lemma~4.6]{Schuch2010PEPS}.
\end{definition}

\begin{theorem}[Coefficient uniqueness for coherent bond sums]%
    \label{thm:peps_four_cut_common_coefficient_uniqueness}
    \lean{TNLean.PEPS.torusBondCoefficientExtraction_bondProduct,
        TNLean.PEPS.torusBondCoefficientExtraction_sum,
        TNLean.PEPS.linearIndependent_torusRepresentationBondProduct,
        TNLean.PEPS.sum_torusRepresentationBondProduct_eq_iff}
    \leanok
    \uses{def:peps_four_cut_common_bond_coefficients}
    Suppose every individual bond representation is semi-regular. Then
    \begin{align}
        E_p(B_q)                      & =\delta_{p,q},
        \label{eq:peps_common_bond_duality}            \\
        E_p\left(\sum_q a_qB_q\right) & =a_p.
        \nonumber
    \end{align}
    The family $(B_p)_p$ is linearly independent. In particular,
    $\sum_pa_pB_p=\sum_pb_pB_p$ if and only if $a_p=b_p$ for every $p$.
\end{theorem}
\begin{proof}\leanok
    \uses{thm:peps_torus_projector_extraction}
    Pair the product vector bond by bond. Trace duality gives
    $E_p(B_q)=\prod_f\delta_{p_f,q_f}=\delta_{p,q}$, since all labelled
    bonds are distinct. Linearity extracts a specified coefficient from
    any finite coherent sum. Applying every $E_p$ to a zero sum proves
    linear independence and applying them to an equality proves coefficient
    uniqueness.
\end{proof}

\begin{theorem}[Relative-label compatibility]%
    \label{thm:peps_four_cut_common_relative_compatibility}
    \lean{TNLean.PEPS.fourBlock_relativeLabel_compatibility,
        TNLean.PEPS.fourBlock_compatibility_commute}
    \leanok
    For arbitrary $a,b,c,d$ in a group, put
    $h_1=b^{-1}a$, $h_2=d^{-1}c$, $v_1=c^{-1}a$, $v_2=d^{-1}b$.
    Then $v_2h_1=h_2v_1$. More generally, any four labels satisfying
    $v_2h_1=h_2v_1$, $h_1=h_2$, and $v_1=v_2$ obey $v_1h_1=h_1v_1$.
    These are the group identities in the proof of
    \cite[Theorem~5.5]{Schuch2010PEPS}.
\end{theorem}
\begin{proof}\leanok
    Both $(d^{-1}b)(b^{-1}a)$ and $(d^{-1}c)(c^{-1}a)$ equal $d^{-1}a$.
    Substituting equal horizontal and vertical labels into the
    compatibility equation gives commutation.
\end{proof}

\begin{definition}[Representation columns and physical bond slices]%
    \label{def:peps_four_cut_common_bond_slices}
    \lean{TNLean.PEPS.representationBondMatrix,
        TNLean.PEPS.representation_matrix_mem_bond_range,
        TNLean.PEPS.torusPhysicalBondSlice,
        TNLean.PEPS.torusBondRepresentation}
    \leanok
    \uses{def:peps_four_cut_common_matched_action,
        def:peps_torus_physical_bond_regrouping}
    Read $U_{(v,\mathrm h)}=U^{\mathrm h}_v$ and
    $U_{(v,\mathrm v)}=U^{\mathrm v}_v$. Define
    $F_f:\C^G\to\C^{V^2}$ by $F_f((i,j),k)=U_f(k)_{i,j}$.
    Thus $\operatorname{ran}F_f$ is the span of the representation matrices,
    and each vector $(i,j)\mapsto U_f(k)_{i,j}$ is its column
    $F_f\delta_k$. For $\tau:\mathcal B\to V^2$, the physical bond slice is
    the linear map
    $S_{f,\tau}\psi(s)=(\mathcal R\psi)(\tau[f\mapsto s])$.
\end{definition}
\begin{proof}\leanok
    Matrix multiplication by the coordinate vector $\delta_k$ gives
    exactly the indicated representation column. Evaluation after the
    inverse bond regrouping is linear in $\psi$.
\end{proof}

\begin{theorem}[An identity bond forces representation-matrix support]%
    \label{thm:peps_four_cut_common_uncut_support}
    \lean{TNLean.PEPS.product_coordinateSlice_eq_smul,
        TNLean.PEPS.torusPhysicalBondSlice_averagingSite_mem_range}
    \leanok
    \uses{def:peps_four_cut_common_bond_slices}
    For any finite set $I$, functions $f_i:Q\to\C$, $j\in I$, and
    $\tau:I\to Q$,
    \begin{align}
        \prod_{i\in I}f_i(\tau[j\mapsto s]_i)
         & =\left(\prod_{i\ne j}f_i(\tau_i)\right)f_j(s).
        \label{eq:peps_common_product_slice}
    \end{align}
    If the inserted bond operator $O_f$ equals $\Id$, then every slice
    $S_{f,\tau}\mathcal N_{A^0}(O)$ lies in $\operatorname{ran}F_f$.
    The other bond operators are arbitrary.
\end{theorem}
\begin{proof}\leanok
    \uses{thm:peps_four_cut_common_network_expansion,
        def:peps_four_cut_common_bond_slices}
    Separate the $j$th factor from the finite product to obtain
    (\ref{eq:peps_common_product_slice}). In
    (\ref{eq:peps_common_regrouped_average}), an identity bond has factor
    $U_f(q_{h(f)}q_{t(f)}^{-1})$. Its $f$-slice is a scalar multiple of
    one column of $F_f$ by (\ref{eq:peps_common_product_slice}). Sum these
    vectors over all $q$ and retain the normalization $|G|^{-|X|}$.
\end{proof}

\begin{theorem}[The four cuts derive a complete bond-product expansion]%
    \label{thm:peps_four_cut_common_bond_span}
    \lean{TNLean.PEPS.torusPhysicalBondSlice_mem_range_of_mem_torusCutSpace,
        TNLean.PEPS.torusBondRegrouping_mem_product_range_of_mem_fourTorusCutSpace,
        TNLean.PEPS.exists_bondCoefficients_of_mem_fourTorusCutSpace,
        TNLean.PEPS.exists_labelCoefficients_of_mem_fourTorusCutSpace,
        TNLean.PEPS.eq_sum_extracted_bondProducts_of_mem_fourTorusCutSpace}
    \leanok
    \uses{def:peps_four_cut_common_boundary,
        def:peps_four_cut_common_bond_slices,
        def:peps_four_cut_common_bond_coefficients,
        def:peps_four_cut_common_physical_product}
    If $\psi\in\mathcal S^{A^0}_{c,r}$ and $f$ is uncut, then
    $S_{f,\tau}\psi\in\operatorname{ran}F_f$ for every $\tau$.
    Here uncut means $v_x+1\ne c$ for $f=(v,\mathrm h)$ or
    $v_y+1\ne r$ for $f=(v,\mathrm v)$.
    On $X=(\Z/2\Z)^2$, every $\psi\in\mathcal S^{A^0}$ satisfies
    \begin{align}
        \mathcal R\psi & \in\operatorname{ran}\mathcal T_F,
        \nonumber                                           \\
        \psi           & =\sum_{q:\mathcal B\to G}a_q
        B_{(v\mapsto q(v,\mathrm h),\,v\mapsto q(v,\mathrm v))}
        =\sum_pa'_pB_p
        \label{eq:peps_common_derived_bond_span}
    \end{align}
    for suitable coefficients. If all bond representations are
    semi-regular, these coefficients are uniquely determined and
    \begin{align}
        \psi & =\sum_pE_p(\psi)B_p.
        \label{eq:peps_common_extracted_expansion}
    \end{align}
\end{theorem}
\begin{proof}\leanok
    \uses{thm:peps_four_cut_common_boundary_span,
        thm:peps_four_cut_common_uncut_support,
        thm:peps_four_cut_common_product_range,
        thm:peps_four_cut_common_coefficient_uniqueness,
        def:peps_four_cut_common_bond_coefficients}
    Expand each boundary in fixed endpoint tensors. Each such cut vector
    has matrix units on cut bonds and identities on uncut bonds, so
    Theorem~\ref{thm:peps_four_cut_common_uncut_support} gives the slice
    condition; linearity extends it to the entire cut range.
    A bond based at $v$ is uncut in the cut $(v_x,v_y)$ on the
    two-by-two torus. Thus every bond satisfies the slice condition, and
    (\ref{eq:peps_common_product_range}) gives a vector $a$ with
    $\mathcal R\psi=\mathcal T_Fa$.
    Expand this product map and undo regrouping to obtain the first sum
    in (\ref{eq:peps_common_derived_bond_span}); the bond-label bijection
    gives the second. Only after this existence argument, apply each
    $E_p$ and use (\ref{eq:peps_common_bond_duality}) to obtain
    (\ref{eq:peps_common_extracted_expansion}).
\end{proof}

\begin{theorem}[Trace extraction of a canonical network]%
    \label{thm:peps_four_cut_common_network_extraction}
    \lean{TNLean.PEPS.torusBondCoefficientExtraction_averagingSite}
    \leanok
    \uses{def:peps_four_cut_common_bond_coefficients}
    For arbitrary inserted bond matrices, with $O^q_f$ as in
    (\ref{eq:peps_common_gauged_horizontal}) and
    (\ref{eq:peps_common_gauged_vertical}),
    \begin{align}
        E_p(\mathcal N_{A^0}(O))
         & =|G|^{-|X|}\sum_{q:X\to G}
        \prod_{f\in\mathcal B}\delta_{f,p_f}(O^q_f).
        \label{eq:peps_common_network_extraction}
    \end{align}
    This identity requires no semi-regularity assumption.
\end{theorem}
\begin{proof}\leanok
    \uses{thm:peps_four_cut_common_network_expansion,
        thm:peps_torus_projector_extraction}
    Apply the linear functional $E_p$ to
    (\ref{eq:peps_common_network_average}). Move it through the finite sum
    and normalization, and factor the pairing of each product into its
    individual bond pairings.
\end{proof}

\begin{theorem}[Nonflat coefficients vanish in the four-cut intersection]%
    \label{thm:peps_four_cut_common_nonflat_vanish}
    \lean{TNLean.PEPS.torusBondCoefficientExtraction_cutNetwork_eq_zero,
        TNLean.PEPS.torusBondCoefficientExtraction_eq_zero_of_mem_torusCutSpace,
        TNLean.PEPS.torusBondCoefficientExtraction_eq_zero_of_mem_fourTorusCutSpace_not_flat}
    \leanok
    \uses{def:peps_four_cut_common_boundary,
        def:peps_four_cut_common_bond_coefficients,def:peps_bond_flat}
    Assume $X=(\Z/2\Z)^2$ and that every bond representation is
    semi-regular. Fix a cut $(c,r)$ and a labelling $p$ for which
    \begin{align}
        p^{\mathrm v}_{(c+1,r)}p^{\mathrm h}_{(c,r+1)}
         & \ne p^{\mathrm h}_{(c,r)}p^{\mathrm v}_{(c,r)}.
        \label{eq:peps_common_failed_plaquette}
    \end{align}
    If $O^{\mathrm h}_v=\Id$ when $v_x+1\ne c$ and
    $O^{\mathrm v}_v=\Id$ when $v_y+1\ne r$, then
    $E_p(\mathcal N_{A^0}(O))=0$.
    Consequently $E_p(\psi)=0$ for every
    $\psi\in\mathcal S^{A^0}_{c,r}$.
    If $\psi\in\mathcal S^{A^0}$ and $p$ is not flat, then
    $E_p(\psi)=0$.
\end{theorem}
\begin{proof}\leanok
    \uses{thm:peps_four_cut_common_network_extraction,
        thm:peps_torus_projector_extraction,
        thm:peps_four_cut_common_boundary_span}
    A nonzero summand in (\ref{eq:peps_common_network_extraction})
    has no zero bond pairing. Trace duality on its uncut bonds forces
    $p^{\mathrm h}_v=q_{v+e}q_v^{-1}$ and
    $p^{\mathrm v}_v=q_vq_{v+n}^{-1}$ there. The four uncut edges then give
    \begin{align}
        p^{\mathrm v}_{(c+1,r)}p^{\mathrm h}_{(c,r+1)}
         & =q_{(c+1,r)}q_{(c,r+1)}^{-1}
        =p^{\mathrm h}_{(c,r)}p^{\mathrm v}_{(c,r)},
        \nonumber
    \end{align}
    contradicting (\ref{eq:peps_common_failed_plaquette}). Thus every
    summand vanishes. Apply this conclusion to the matrix-unit networks
    spanning the cut range. Finally, a nonflat labelling fails the
    plaquette equation at some $(c,r)$, whose cut membership gives the result.
\end{proof}

\begin{definition}[Product of the local invariant projectors]%
    \label{def:peps_four_cut_common_average_projection}
    \lean{TNLean.PEPS.representationAveragingSite_invariant,
        TNLean.PEPS.torusSitewiseAveragingProjector}
    \leanok
    \uses{def:peps_four_cut_common_matched_action,
        def:peps_four_cut_common_physical_product}
    For arbitrary representations $\rho_v$ of finite $G$ on $\C^{V^4}$,
    define $\mathcal Q_\rho=\mathcal T_{(\Pi_{\rho_v})_v}$.
    Each averaging tensor is invariant:
    $\mathcal P(A^0_{\rho_v})\rho_v(k)=\mathcal P(A^0_{\rho_v})$ for all $k$.
\end{definition}
\begin{proof}\leanok
    Right multiplication by $k$ permutes the summands of
    $\Pi_{\rho_v}=|G|^{-1}\sum_g\rho_v(g)$, so
    $\Pi_{\rho_v}\rho_v(k)=\Pi_{\rho_v}$.
\end{proof}

\begin{theorem}[Projection fixes cut vectors and inserts averaging sites]%
    \label{thm:peps_four_cut_common_average_projection}
    \lean{TNLean.PEPS.torusSitewiseAveragingProjector_cutMap,
        TNLean.PEPS.torusSitewiseAveragingProjector_eq_self_of_mem_cut,
        TNLean.PEPS.torusSitewiseAveragingProjector_bondProduct}
    \leanok
    \uses{def:peps_four_cut_common_average_projection,
        def:peps_four_cut_common_boundary,def:peps_four_cut_common_bond_coefficients}
    For every boundary $M$,
    $\mathcal Q_\rho\mathcal C^{A^0}_{c,r}M=\mathcal C^{A^0}_{c,r}M$.
    Therefore $\mathcal Q_\rho\psi=\psi$ for every
    $\psi\in\mathcal S^{A^0}_{c,r}$. For any bond representations
    $U^{\mathrm h},U^{\mathrm v}$ and labelling $p$,
    \begin{align}
        \mathcal Q_\rho B_p
         & =\mathcal N_{A^0}
        ((U^{\mathrm h}_v(p^{\mathrm h}_v))_v,
        (U^{\mathrm v}_v(p^{\mathrm v}_v))_v).
        \label{eq:peps_common_project_bond_product}
    \end{align}
    The actions $\rho_v$ need not be the matched actions for these bond
    representations.
\end{theorem}
\begin{proof}\leanok
    \uses{thm:peps_four_cut_common_local_inverse,
        def:peps_four_cut_common_average_projection,
        thm:peps_four_cut_common_physical_contraction,
        thm:peps_torus_projector_expansion}
    Apply the recovery identity
    (\ref{eq:peps_common_cut_recovery}) with $A_v=A^0_{\rho_v}$;
    its invariance proves the fixed-vector assertion.
    Express $B_p$ as a bond network whose local tensors have coefficients
    $\delta_{(t,r,b,l),s}$. Applying $\Pi_{\rho_v}$ to such an identity
    tensor produces $A^0_{\rho_v}$. Physical transport through the
    contraction gives (\ref{eq:peps_common_project_bond_product}).
\end{proof}

\begin{theorem}[Projecting a flat product gives a commuting closure]%
    \label{thm:peps_four_cut_common_flat_projection}
    \lean{TNLean.PEPS.exists_averagingProjector_bondProduct_eq_matchedClosure,
        TNLean.PEPS.averagingProjector_bondProduct_mem_matchedClosureSpan}
    \leanok
    \uses{def:peps_four_cut_common_average_projection,
        def:peps_four_cut_common_bond_coefficients,
        def:peps_four_cut_common_closure,def:peps_bond_flat}
    Use the matched actions $\rho_v$. For every flat bond labelling $p$
    on a torus with positive periods, there exist $g,h\in G$ with $gh=hg$
    such that $\mathcal Q_\rho B_p=\mathcal M_{A^0}(g,h)$.
    In particular, $\mathcal Q_\rho B_p\in\mathcal Z_{A^0}$.
\end{theorem}
\begin{proof}\leanok
    \uses{thm:peps_bond_flat_classification,
        thm:peps_four_cut_common_average_projection,
        thm:peps_four_cut_common_vertex_gauge,
        def:peps_four_cut_common_average_projection}
    Theorem~\ref{thm:peps_bond_flat_classification} gives a vertex gauge
    $q$ sending $p$ to two seams labelled by commuting $g,h$.
    Formula (\ref{eq:peps_common_project_bond_product}) expresses
    $\mathcal Q_\rho B_p$ as the corresponding actual averaging network.
    Its invariant sites permit the matched vertex-gauge transformation.
    Applying each bond's representation to the transformed labels turns
    its matrices into $H^0(h)$ and $V^0(g)$, giving
    $\mathcal M_{A^0}(g,h)$. This vector is a generator of the stated span.
\end{proof}

\begin{theorem}[Canonical four-cut closure with common alphabets]%
    \label{thm:peps_four_cut_common_canonical_closure}
    \lean{TNLean.PEPS.fourTorusCutSpace_matchedAveraging_le_commutingClosureSpan,
        TNLean.PEPS.fourTorusCutSpace_matchedAveraging_eq_commutingClosureSpan}
    \leanok
    \uses{def:peps_four_cut_common_boundary,def:peps_four_cut_common_closure}
    On the two-by-two torus, assume every individual bond representation is
    semi-regular. Then $\mathcal S^{A^0}\subseteq\mathcal Z_{A^0}$ and,
    together with the opposite inclusion,
    $\mathcal S^{A^0}=\mathcal Z_{A^0}$.
\end{theorem}
\begin{proof}\leanok
    \uses{thm:peps_four_cut_common_bond_span,
        thm:peps_four_cut_common_nonflat_vanish,
        thm:peps_four_cut_common_average_projection,
        thm:peps_four_cut_common_flat_projection,
        thm:peps_four_cut_common_matched_membership,
        def:peps_four_cut_common_average_projection}
    For $\psi\in\mathcal S^{A^0}$, the cut $(0,0)$ gives
    $\mathcal Q_\rho\psi=\psi$. Apply this projector to the derived
    expansion (\ref{eq:peps_common_extracted_expansion}):
    \begin{align}
        \psi & =\sum_pE_p(\psi)\mathcal Q_\rho B_p
        =\sum_{p\text{ flat}}E_p(\psi)\mathcal Q_\rho B_p
        \in\mathcal Z_{A^0}.
        \label{eq:peps_common_flat_reconstruction}
    \end{align}
    The nonflat terms have zero extracted coefficient, while every flat
    term is a commuting closure. The averaging sites are invariant, so
    Theorem~\ref{thm:peps_four_cut_common_matched_membership} gives
    $\mathcal Z_{A^0}\subseteq\mathcal S^{A^0}$.
\end{proof}

\begin{theorem}[Four-cut closure for independent tensors with common alphabets]%
    \label{thm:peps_four_cut_common_closure}
    \lean{TNLean.PEPS.fourTorusCutSpace_eq_matchedCommutingClosureSpan}
    \leanok
    \uses{def:peps_four_cut_common_boundary,
        def:peps_four_cut_common_matched_action,
        def:peps_four_cut_common_closure,def:peps_g_injective}
    Let $G,V,P$ be finite and $X=(\Z/2\Z)^2$. Assign an arbitrary
    semi-regular representation on $\C^V$ to each of the eight bonds,
    and let each independent site tensor $A_v:V^4\times P\to\C$ be
    $G$-injective for its matched action $\rho_v$. Then
    \begin{align}
        \bigcap_{c,r\in\Z/2\Z}\operatorname{ran}\mathcal C^A_{c,r}
         & =\spn_{\C}\{\mathcal M_A(g,h):g,h\in G,\ gh=hg\}.
        \label{eq:peps_common_four_cut_closure}
    \end{align}
    This is the common-coordinate specialization of
    \cite[Theorem~5.5]{Schuch2010PEPS}; the independently sized bond and
    physical spaces are treated in Theorem~\ref{thm:peps_four_cut_closure}.
\end{theorem}
\begin{proof}\leanok
    \uses{thm:peps_four_cut_common_local_inverse,
        thm:peps_four_cut_common_canonical_closure}
    Put $R_v=\mathcal P(A_v)$. The common local inverse and the canonical
    equality give
    \begin{align}
        \mathcal S^A
         & =\mathcal T_R(\mathcal S^{A^0})
        =\mathcal T_R(\mathcal Z_{A^0})
        =\mathcal Z_A.
        \nonumber
    \end{align}
    The last equality is closure recovery by the original invariant site
    maps. Each cut retains its arbitrary joint boundary throughout.
\end{proof}
